\section{What accumulates before the cut-off and what remains at the boundary}\label{sec:signedmap}
Choose a solution uniformly and reveal its queens in a uniformly random row order independent of the
solution. Revealing the rows in random order is the device of the entropy method~\cite{rad,ll}, which
gives the upper bounds on \(Q(n)\) in~\cite{luria,simkin}; here the same histories are followed under
two laws. When \(m\) rows remain, the residual board \(B_m\) records the available squares. Its binary
mask \(M_m\) supports a simpler problem: place one rook in each row and column. The number of
rook placements, \(\operatorname{per}M_m\), already incorporates the restrictions imposed by the revealed
pieces. To obtain \(Q(B_m)\) we must also require that the new pieces do not attack each other along the
diagonals.

This suggests correcting the queen predictor by the exact rook count. Let \(L\) and \(L_T\) be the
logarithms of the queen and rook predictors on
the residual board. For the rook part we can use the exact count:
\[
 g_T(M)=\log\operatorname{per}M-L_T(M),\qquad L^\sharp(B)=L(B)+g_T(M).
\]
The adjustment is made on the same mask at each step. Thus a difference between two child boards does not
mix the error in counting rooks with the cost of forbidding a diagonal attack. On the initial board \(n
g_T(M_n)\) tends to \(F(1)-1/12\). The full mask is separable and \(F(1)=0\); this gives the first
summand in~\eqref{eq:main}.

\paragraph{The sum up to the stopping time.} We stop when \(j=\lfloor bn\rfloor\) empty rows remain, or
earlier at the first exit. The proof distinguishes two levels. The first, \(M=\lceil\delta n\rceil\),
separates the regularity regimes of~\cite[Definition~2.2 and \S5.4]{tc}. Until \(M\) we require the
discrepancy of the history from the profile to stay below its radius. At \(M\) we also check the entry
tolerances and the maximum diagonal length. From \(M\) to \(j\) we require the inner margins for
diffusion of the weights, vacancy, spectral gap, and diagonal mass. The first board that violates any
of these conditions counts as an exit, even if it is reached precisely at \(j\).

With \(0<\delta<1\) and these margins fixed, we choose a length threshold \(\alpha>0\) and then an
admissible cut-off \(b>0\). The restrictions on \(b\) are specified in Appendix~\ref{app:arrival}; they
allow all sufficiently small \(b\) with \(b<\delta/2\). Only then do we take \(n\) large. We write
\(F_b\) for the event of reaching \(j\) without an exit and with at least one completion. The \emph{short
historical arrival}, \(H_j\), requires regular entry at \(M\), persistence until \(j\), and fewer than
\(\lceil\alpha j\rceil\) available squares on each diagonal and antidiagonal at arrival. By
Lemma~\ref{lem:positividad-llegada}, on \(H_j\) the mask admits a queen completion, so \(H_j\subseteq
F_b\). In each row the adaptive law chooses a square with probability proportional to the weight of its
child board. If the sum of these weights equaled the weight of the parent, the recursion would hold
exactly in that row. The logarithm of their ratio measures the defect of that step. We denote the sum of
these logarithms up to the stopping time by \(A^\sharp\).

The uniform solution and the adaptive law use the same stopping rule. This lets us compare the
probabilities of the same history and retain the sign of the difference in exit probabilities.
Appendix~\ref{app:signed} gives the exact identity, and Appendix~\ref{app:arrival} controls the remainder
in the mean. At fixed \(b\) the defect accumulated in the interior contributes \(\int_b^1
\numerr(x)x^{-2}\,dx\). Appendix~\ref{app:interior} identifies its sixteen terms.

\paragraph{What survives on the arrival board.} Consider two available squares \(e,f\) on the same
diagonal. They are in distinct rows and columns, so two rooks could occupy them simultaneously. The other
two positions of the square determine four states. For \(e=(r,q)\) and \(f=(r+h,q+h)\), with \(h>0\), let
\(X\) and \(Y\) be the availability indicators of \((r,q+h)\) and \((r+h,q)\), respectively.
Figure~\ref{fig:frontera} shows the four cases.
\begin{figure}[htbp]
\centering
\begin{tikzpicture}[x=.58cm,y=.58cm]
 \foreach \dx/\sx/\sy in {0/0/0,5.8/1/0,11.6/0/1,17.4/1/1}{
  \begin{scope}[shift={(\dx,0)}]
   \bcuadricula{4}
   \draw[btinta,line width=.9pt] (1.5,3.5)--(3.5,1.5);
   \ifnum\numexpr\sx+\sy\relax=2\relax
    \draw[btinta,dash pattern=on 2pt off 1.5pt,line width=.9pt] (1.5,1.5)--(3.5,3.5);
   \fi
   \foreach \fr/\fc/\libre in {1/2/1,3/4/1,1/4/\sx,3/2/\sy}{
    \begin{scope}[shift={(\fc-.5,4.5-\fr)}]
     \fill[bpapel] (0,0) circle[radius=.26];
     \ifnum\libre=1\relax
      \draw[btinta,line width=1pt] (0,0) circle[radius=.22];
     \else
      \draw[btinta,line width=1.05pt] (-.17,-.17)--(.17,.17) (-.17,.17)--(.17,-.17);
     \fi
    \end{scope}
    \ifnum\libre=1\relax \bmarca{4}{\fr}{\fc}\fi
   }
   \node[anchor=north,font=\small] at (2,-.18) {$(X,Y)=(\sx,\sy)$};
  \end{scope}
 }
\end{tikzpicture}
\caption{The four states of the boundary square are drawn in a \(4\times4\) window of a larger board,
with local coordinates \(r=1\), \(q=2\), \(h=2\). A circle indicates an available square and a cross an
excluded one. The solid line joins \(e=(1,2)\) and \(f=(3,4)\). The dashed line joins \((1,4)\) and
\((3,2)\) when both are available. The other squares are shown only for reference. In state \(11\), after
subtracting the effects of the individual pairs, the interaction coefficients in the logarithms of the
count and of the predictor are \(-1\) and \(+2\), in units of the product of the four weights at leading
order. Their difference \(-1-(+2)=-3\) gives the sign of \(-S\) in~\eqref{eq:exceso-cuadrado}.}
\label{fig:frontera}
\end{figure}


All four cases occur in the solution \((1,3,5,7,2,4,6)\) on the board of order seven, for the pair
\(e=(1,2)\), \(f=(3,4)\) of Figure~\ref{fig:frontera}, whose other two vertices are \((1,4)\) and
\((3,2)\). Before any queen is
revealed, \(X=Y=1\). Revealing \((7,6)\) rules out only \(Y\) among the four vertices, and revealing
\((4,7)\) rules out only \(X\). If both are revealed, \(e,f\) remain available.

The response of the marked pair is evaluated with \(X,Y\) fixed. In state \(11\) either diagonal pair may
be marked, and each contributes its own response. Both belong to the same square, so their interaction
enters the sum over complete squares once. We denote the sum of the three boundary means by \(\mathcal
M_n(b)\). Letting \(n\to\infty\) and then \(b\downarrow0\), we obtain \(P_T\) and \(J\) from the two means
over pairs and \(-S\) from the mean interaction. Appendices~\ref{app:boundary} and~\ref{app:weights}
separate the contributions. Appendix~\ref{app:transfer} averages the pair response under the arrival law.

\paragraph{Why compare two cut-offs?} For \(\sumafrontera=P_T+J-S\), write
\[
 \valorcorte(b)=F(1)-\tfrac1{12}+\int_b^1 \numerr(x)x^{-2}\,dx+\sumafrontera.
\]
The error satisfies \(\limsup_n |nR_Q(n)-\valorcorte(b)|\le C_\delta b+\eta(b)\), where
\(\eta(b)=\limsup_n|\mathcal M_n(b)-\sumafrontera|\to0\) as \(b\) decreases. The constant \(C_\delta\)
works for all sufficiently small admissible cut-offs. Each choice of \(b\) approximates the same sequence
\(nR_Q(n)\). As \(b\) and \(b'\) tend to zero, the difference between the two approximations tends to
zero. This yields a common limit before the integral is extended to the endpoint. Once poles are ruled out
as in Lemma~\ref{lem:regularidad-extremo}, the constant term of \(\numerr\) would produce a divergence of
order \(1/b\) in the truncated integral and the linear term one of order \(\log(1/b)\). The existence of
a finite limit rules out both, in that order (Appendix~\ref{app:closing}).

Figure~\ref{fig:cortes} illustrates the two levels in a particular history. In the proof, each admissible
cut-off gives an approximation through averages over histories. Since both approximations refer to the
same \(nR_Q(n)\), the triangle inequality gives \(\lvert\valorcorte(b)-\valorcorte(b')\rvert\le
C_\delta(b+b')+\eta(b)+\eta(b')\).

\begin{figure}[htbp]
\centering
\begin{tikzpicture}[x=.58cm,y=.58cm]
 \foreach \etapa/\dx in {2/0,4/10.4}{
  \begin{scope}[shift={(\dx,0)}]
   \bcuadricula{6}
   \breina{6}{1}{2}
   \breina{6}{2}{4}
   \ifnum\etapa=2\relax
    \def\casillas{3/1,3/6,4/1,4/3,5/3,5/5,6/1,6/3,6/5,6/6}
    \def\corte{b=2/3}
   \else
    \breina{6}{3}{6}
    \breina{6}{4}{1}
    \def\casillas{5/3,5/5,6/5}
    \def\corte{b'=1/3}
   \fi
   \foreach \fila/\col in \casillas{
    \draw[btinta,line width=.55pt] (\col-.5,6.5-\fila) circle[radius=.16];
    \bmarca{6}{\fila}{\col}
   }
   \draw[btinta,line width=.95pt] (0,6-\etapa)--(6,6-\etapa);
   \node[font=\small,anchor=north] at (3,-.18) {$\corte,\quad m=\the\numexpr6-\etapa\relax$};
  \end{scope}
 }
 \draw[->,>=stealth,btinta,line width=.8pt] (6.55,3)--(9.45,3)
   node[midway,above,font=\scriptsize] {rows 3 and 4};
\end{tikzpicture}
\caption{Two cut-offs of the solution \((2,4,6,1,3,5)\). Entry \(r\) gives the column of the queen in row
\(r\). In this example the rows are revealed from top to bottom, and the horizontal line marks how far
the process has gone. Crowns indicate revealed queens, and circles with a central dot indicate the
squares those
queens leave available. After rows \(3\) and \(4\) are revealed, three of the ten remain.}
\label{fig:cortes}
\end{figure}

